\begin{explanationbox}
\textbf{\textcolor{CExplanation}{一句话直觉}}
一个关于三个实数 $a, b, c$ 的“强度”不等式链：最弱的是两两相乘的和，最强的是各自的平方和，中间是它们和的平方的三分之一。

\medskip
\textbf{\textcolor{CExplanation}{核心拆解}}
\begin{description}[style=nextline, leftmargin=2.8em, labelsep=0.5em, itemsep=0.45em]
    \item[\textbf{要点一（链式结构）：}] 它不是一个孤立不等式，而是两个不等式构成的链条，揭示了三个量 $S_1 = ab+bc+ca$， $S_2 = (a+b+c)^2/3$， $S_3 = a^2+b^2+c^2$ 之间固定的强弱关系：$S_1 \le S_2 \le S_3$。
    \item[\textbf{要点二（对称性）：}] 不等式关于 $a, b, c$ 完全对称，任意交换变量，不等式不变。这提示证明时通常利用对称性进行配方。
\end{description}

\medskip
\textbf{\textcolor{CExplanation}{几何本质（平方距离）}}
不等式 $a^2+b^2+c^2 \ge ab+bc+ca$ 等价于 $\frac12[(a-b)^2+(b-c)^2+(c-a)^2] \ge 0$。其几何意义是：在三维空间中，点 $(a,b,c)$ 到直线 $x=y=z$（即“相等线”）的平方距离非负。左边较弱的不等式 $ab+bc+ca \le (a+b+c)^2/3$ 则反映了“乘积和”与“算术平均平方”的关系。

\medskip
\textbf{\textcolor{CExplanation}{代数意义（均值联系）}}
不等式链将三种常见的对称式联系了起来：
\begin{itemize}[leftmargin=*, itemsep=0.3em, topsep=0.2em]
    \item $ab+bc+ca$：两两乘积的初级对称式。
    \item $\dfrac{(a+b+c)^2}{3}$：和的平方与常数 $3$ 的商，也等于 $a,b,c$ 算术平均值的平方的 $3$ 倍。
    \item $a^2+b^2+c^2$：平方和，是二次型的基本形式。
\end{itemize}
链表明，对于任意实数组，初级对称式不超过“平均的平方尺度”，而后者又不超过平方和。

\medskip
\textbf{\textcolor{CExplanation}{考点价值（放缩桥梁）}}
\begin{itemize}[leftmargin=*, itemsep=0.5em, topsep=0.4em]
    \item \textbf{考法一（最值估算）：} 已知 $a+b+c$ 的值，可立刻估计 $ab+bc+ca$ 的最大值，或已知 $a^2+b^2+c^2$ 的值，估计 $a+b+c$ 的范围。
    \item \textbf{考法二（不等式证明）：} 在复杂不等式证明中，常用此链作为放缩的桥梁，将难以处理的 $ab+bc+ca$ 转化为平方和或和的平方进行处理。
\end{itemize}

\medskip
\textbf{\textcolor{CExplanation}{顿悟点}}
\textbf{这个链条是“平方非负”这一最基本事实的集中体现。} 右侧 $a^2+b^2+c^2 \ge ab+bc+ca$ 直接由 $(a-b)^2+(b-c)^2+(c-a)^2 \ge 0$ 推出；左侧 $ab+bc+ca \le (a+b+c)^2/3$ 也由 $(a+b+c)^2 - 3(ab+bc+ca) \ge 0$ 推出，而后者同样是平方和形式。所以，它的本质是“配方”。

\medskip
\textbf{\textcolor{CExplanation}{使用场景}}
\begin{itemize}[leftmargin=*, itemsep=0.5em, topsep=0.4em]
    \item \textbf{场景一（条件最值）：} 当题目条件给出 $a+b+c$ 为定值，求 $ab+bc+ca$ 的最大值时，直接使用左半部分不等式。
    \item \textbf{场景二（对称放缩）：} 在证明含 $a^2+b^2+c^2$、$ab+bc+ca$、$(a+b+c)^2$ 的对称不等式时，可尝试用此链进行统一或转化。
\end{itemize}
\end{explanationbox}
